% !TEX root = ../main.tex
% ══════════════════════════════════════════════════════════════
\chapter{Positionnement par rapport aux solutions existantes}
% ══════════════════════════════════════════════════════════════
\begingroup
\let\clearpage\relax
\startcontents[chapters]
\printcontents[chapters]{}{1}{\section*{Plan}}
\endgroup
\newpage

\section*{Introduction}
\justifying
\phantomsection
\addcontentsline{toc}{section}{Introduction}

Le chapitre précédent a mis en évidence les limites du système d'information du service RI et formulé la problématique du projet. Ce chapitre évalue les solutions existantes au regard des huit critères qui en découlent, afin de vérifier qu'aucune ne répond déjà à ces besoins avant d'engager le développement d'une solution sur mesure.

% ─────────────────────────────────────────────────────────────
\section{Grille d'évaluation}
% ─────────────────────────────────────────────────────────────

Pour comparer les solutions de façon rigoureuse, une grille de huit critères a été définie
à partir des problèmes concrets identifiés au chapitre précédent. Chaque critère traduit une
lacune précise du système actuel. Le tableau~\ref{tab:grille_eval} présente cette grille.

\renewcommand{\arraystretch}{1.45}
\begin{table}[H]
\centering
\footnotesize
\begin{tabular}{|p{0.6cm}|p{3.8cm}|p{9.2cm}|}
\hline
\rowcolor{gray!20}
\textbf{\#} & \textbf{Critère} & \textbf{Ce que cela signifie concrètement pour l'ENSEEIHT} \\
\hline
C1 & Centralisation des données de mobilité & Regrouper dans une seule plateforme
l'ensemble des données relatives aux mobilités sortantes, entrantes, complémentaires
et aux stages internationaux, ainsi que les données académiques des étudiants \\
\hline
C2 & Quotas par école et par département & Gestion des places allouées d'abord à l'échelle de l'ENSEEIHT, puis par département pédagogique (SN, 3EA, MF2E), et non uniquement à l'échelle globale INP \\
\hline
C3 & Aide à la décision et auto-affectation & Présence d'un mécanisme automatique de classement et de pré-affectation des candidats. \\
\hline
C4 & Calcul de la durée de mobilité & Agrégation automatique de toutes les périodes de mobilité d'un étudiant (académique, stage, complémentaire) pour produire la durée réglementaire \\
\hline
C5 & Mobilités complémentaires & Prise en charge des expériences hors accord bilatéral (summer schools, séjours linguistiques, programmes courts) avec dépôt et validation de justificatifs \\
\hline
C6 & Adaptabilité aux règles internes & Capacité du système à fonctionner selon les règles propres à l'établissement : priorités d'affectation, seuils de sélection, critères locaux \\
\hline
C7 & Recommandation de destinations & Suggestion de destinations adaptées au profil de l'étudiant (département, GPA, langue, préférences) pour l'accompagner dans son choix \\
\hline
C8 & Reporting CTI & Production automatique des indicateurs clés (KPIs) exigés par la CTI, notamment : taux de mobilité par promotion, durée moyenne de mobilité par étudiant, et répartition géographique des destinations \\
\hline
\end{tabular}
\caption{Grille d'évaluation des solutions existantes}
\label{tab:grille_eval}
\end{table}

% ─────────────────────────────────────────────────────────────
\section{Analyse des solutions existantes}
% ─────────────────────────────────────────────────────────────

Quatre solutions ont été examinées : QS MoveON, déjà utilisé par le service RI, Mobility-Online, une plateforme européenne largement déployée dans les services de relations internationales, ErasmusJET, développé dans le cadre du programme Erasmus+, et les travaux académiques d'Elahi et al. sur la recommandation de destinations.

\subsection{QS MoveON}

MoveON est l'outil central du service RI pour la gestion des accords bilatéraux et des candidatures. Il est examiné ici non pas pour le remplacer, mais pour identifier précisément ce qu'il ne couvre pas~\cite{qs_moveon}.

\textbf{Points forts  :} le référentiel des accords bilatéraux, le suivi des candidatures post-sélection et l'administration des dossiers étudiants. Son API REST est un point d'appui utile pour extraire les données.

\textbf{Limites par rapport aux critères définis :}

\begin{itemize}
    \item MoveON se limite aux mobilités sortantes et aux accords. Les données académiques (GPA, département, parcours), les stages internationaux et les mobilités complémentaires en sont totalement absents. Le service RI est donc contraint d'exporter depuis MoveON, de compléter à la main dans Excel, puis de travailler sur ce fichier hybride — ce cumul de gestes est exactement ce que ce projet cherche à éliminer.

    \item Les quotas ne sont gérés qu'à l'échelle INP. La déclinaison par école (ENSEEIHT) puis par département pédagogique, qui conditionne l'affectation interne, est entièrement tenue dans des fichiers Excel sans aucune contrainte de cohérence.

    \item Aucun classement automatique n'existe. La sélection est intégralement manuelle, ce qui signifie que les règles appliquées varient selon l'agent et la campagne.

    \item Sans les stages ni les expériences complémentaires, le calcul de la durée de mobilité est impossible à automatiser depuis MoveON seul.

    \item Aucune suggestion de destinations adaptée au profil de l'étudiant.

    \item Les exports MoveON ne produisent pas les indicateurs CTI structurés (taux de mobilité par promotion, durées moyennes, répartition géographique) dans un format directement exploitable pour les audits.
\end{itemize}

MoveON reste un bon outil pour ce qu'il fait, mais sa conception ne prévoit pas de
couvrir les besoins opérationnels internes d'un établissement.

\subsection{Mobility-Online}

Mobility-Online est une plateforme de gestion de la mobilité étudiante largement utilisée par des services de relations internationales en Europe, notamment pour la gestion administrative des échanges Erasmus+ et des accords bilatéraux.

\textbf{Points forts :} un module de gestion documentaire complet (conventions, learning agreements, relevés de notes) et une bonne couverture du suivi administratif des mobilités entrantes et sortantes, sur un périmètre européen large.

\textbf{Limites par rapport aux critères définis :}

\begin{itemize}
    \item La plateforme est conçue comme un outil de gestion administrative générique, sans granularité de quotas par département pédagogique propre à un établissement français comme l'ENSEEIHT.
    \item Aucun mécanisme d'auto-affectation basé sur le classement académique n'est proposé : la sélection reste manuelle.
    \item Le calcul de durée de mobilité selon les critères CTI, spécifiques aux écoles d'ingénieurs françaises, n'est pas pris en charge, la plateforme n'étant pas conçue pour ce cadre réglementaire.
    \item Ni recommandation de destinations personnalisée, ni reporting structuré selon les indicateurs CTI.
\end{itemize}

Mobility-Online confirme qu'une plateforme généraliste, même mature et largement déployée, ne couvre pas les besoins spécifiques d'un établissement soumis aux exigences de la CTI.

\subsection{ErasmusJET}

ErasmusJET est une plateforme développée dans le cadre des initiatives de digitalisation du programme Erasmus+. Son objectif est de standardiser les échanges de données entre établissements partenaires du réseau européen.

\textbf{Points forts :} la gestion complète du cycle Erasmus — dépôt de candidatures en ligne, génération et signature des conventions, échange de relevés de notes entre institutions. Pour un établissement dont toutes les mobilités relèvent du programme Erasmus, c'est une solution adaptée.

\textbf{Limites par rapport aux critères définis :}

\begin{itemize}
    \item ErasmusJET ne fonctionne que pour Erasmus+. Les accords hors Europe — qui représentent une grande partie des mobilités de l'ENSEEIHT — restent en dehors de son périmètre, tout comme les données académiques, les stages et les mobilités complémentaires.

    \item Les quotas par département, l'auto-affectation et le calcul de durée CTI ne sont pas prévus dans la plateforme.

    \item Son format est standardisé pour les échanges européens, ce qui ne laisse pas de place pour les règles propres à chaque établissement.

    \item Ni recommandation de destinations, ni reporting CTI.
\end{itemize}

ErasmusJET peut être utile pour la partie Erasmus+, mais il ne règle aucun des problèmes identifiés.

\subsection{Travaux académiques : Elahi et al. (2022)}

Les travaux d'Elahi et al. proposent un système qui combine recommandation de destinations et processus d'affectation étudiante. Les tests sont réalisés sur un jeu de données simulé~\cite{Elahi2022}.

\textbf{Points forts :} c'est la seule des solutions analysées qui aborde sérieusement le critère C7 — la recommandation personnalisée de destinations. L'article démontre la faisabilité technique de coupler un algorithme de suggestion avec un processus de sélection, ce qui en fait une référence directement pertinente pour la partie intelligence du projet.

\textbf{Limites par rapport aux critères définis :}

\begin{itemize}
    \item Le modèle repose sur un dataset simulé et uniforme, sans connexion avec des systèmes institutionnels réels. Il ne propose aucune intégration avec MoveON, Pégase ou Eudonet.

    \item La gestion des quotas par département, l'affectation sous contraintes, le calcul de durée de mobilité et les expériences complémentaires sont entièrement hors du périmètre — le modèle ne traite que la recommandation.

    \item C'est un prototype de recherche : il n'y a aucun moyen de configurer des règles propres à un établissement.

    \item Le reporting CTI n'est pas abordé.
\end{itemize}

Ces travaux sont une référence utile pour la partie recommandation, mais ils ne forment pas une solution prête à être utilisée à l'ENSEEIHT.

% ─────────────────────────────────────────────────────────────
\section{Tableau comparatif de synthèse}
% ─────────────────────────────────────────────────────────────

Le tableau~\ref{tab:synthese_solutions} résume la couverture de chaque solution pour les huit critères de la grille d'évaluation. La légende suivante s'applique : \cmark\ = critère satisfait, \xmark\ = critère non satisfait, \textit{Partiel} = critère partiellement couvert.

\renewcommand{\arraystretch}{1.45}
\begin{table}[H]
\centering
\footnotesize
\begin{tabular}{|p{5.4cm}|>{\centering\arraybackslash}p{2cm}|>{\centering\arraybackslash}p{2.1cm}|>{\centering\arraybackslash}p{2cm}|>{\centering\arraybackslash}p{2.1cm}|}
\hline
\rowcolor{gray!20}
\centering{\textbf{Critère}} & \textbf{QS MoveON} & \textbf{Mobility-Online} & \textbf{ErasmusJET} & \textbf{Elahi et al.} \\
\hline
C1 – Centralisation des données & Partiel & Partiel & \xmark & \xmark \\
\hline
C2 – Quotas par école et département & \xmark & \xmark & \xmark & \xmark \\
\hline
C3 – Aide à la décision et auto-affectation & \xmark & \xmark & \xmark & \xmark \\
\hline
C4 – Calcul durée de mobilité & \xmark & \xmark & \xmark & \xmark \\
\hline
C5 – Mobilités complémentaires & \xmark & \xmark & \xmark & \xmark \\
\hline
C6 – Adaptabilité métier & Partiel & \xmark & \xmark & \xmark \\
\hline
C7 – Recommandation de destinations & \xmark & \xmark & \xmark & Partiel \\
\hline
C8 – Reporting CTI & \xmark & \xmark & \xmark & \xmark \\
\hline
\end{tabular}
\caption{Synthèse comparative des solutions selon la grille d'évaluation}
\label{tab:synthese_solutions}
\end{table}

Aucune des quatre solutions étudiées — qu'il s'agisse d'outils commerciaux largement déployés (MoveON, Mobility-Online), d'une initiative institutionnelle européenne (ErasmusJET) ou d'un prototype de recherche (Elahi et al.) — ne répond à l'ensemble des besoins identifiés. Ce résultat, obtenu sur un panel volontairement diversifié, confirme la nécessité de développer une solution propre au service RI, conformément à la problématique et à la solution envisagées au chapitre précédent.

\section*{Conclusion}
\justifying
\phantomsection
\addcontentsline{toc}{section}{Conclusion}

Ce chapitre a montré qu'aucune des solutions existantes — commerciales ou académiques — ne couvre les huit critères identifiés, en particulier la granularité des quotas par département et la conformité CTI, propres au contexte français. Ce constat valide la solution envisagée au chapitre précédent ; le chapitre suivant en détaille les choix technologiques et algorithmiques.
